%! Matrix Computations and A = LU — Unit 2, Lesson 2.3 — Group Activity (Answer Key)
\documentclass[10pt]{article}
\usepackage{linalg-article}
\usepackage{linalg-key}   % pulls in -boxes; do NOT also load -boxes
\newcommand{\TallMath}[1]{$\displaystyle #1\rule[-1.4em]{0pt}{3.2em}$}
\newcommand{\xx}{\mathbf{x}}
\newcommand{\bb}{\mathbf{b}}
\newcommand{\cc}{\mathbf{c}}

\begin{document}

\pageheader{Unit 2, Lesson 2.3}{Group Activity --- Answer Key}
\namepartnerperiod

\vspace{0.10in}

\begin{headlinebox}{blush}
\small\textbf{Factor Once, Solve Fast.} A factorization $A=LU$ \emph{saves} elimination: $U$ is the
upper-triangular result, and $L$ is lower-triangular with $1$s on the diagonal and the multipliers
$\ell$ below. To solve $A\xx=\bb$, sweep twice --- forward $L\cc=\bb$, then back $U\xx=\cc$. Work with
your partner and \emph{show every step}.
\end{headlinebox}

\vspace{0.10in}

\begin{tcolorbox}[colback=white, colframe=black!40, title=\textbf{Tier R --- Remediate},
    fonttitle=\bfseries, arc=1.5mm, left=3mm, right=3mm, top=2mm, bottom=2mm, breakable]
Factor $A=\begin{bmatrix} 2 & 6 \\ 6 & 20 \end{bmatrix}$. The pivot is $2$, so the multiplier is
$\ell=6\div2=3$.
\begin{enumerate}[leftmargin=1.7em, itemsep=7pt]
  \item Eliminate (row$_2 - 3\,$row$_1$) to get $U$. \; $U=\begin{bmatrix} 2 & 6 \\ {\color{keyred}\mathbf{0}} & {\color{keyred}\mathbf{2}} \end{bmatrix}$
        \hfill {\footnotesize$-3(2,6)+(6,20)=(0,2)$}
  \item Write $L$: put $1$s on the diagonal and the multiplier below. \; $L=\begin{bmatrix} 1 & 0 \\ {\color{keyred}\mathbf{3}} & 1 \end{bmatrix}$
  \item \textbf{Check} $LU=A$ by multiplying.
        \ansline{$\begin{bsmallmatrix} 1&0\\3&1 \end{bsmallmatrix}\begin{bsmallmatrix} 2&6\\0&2 \end{bsmallmatrix}=\begin{bsmallmatrix} 2&6\\6&20 \end{bsmallmatrix}=A$ \checkmark \ (row~2: $3(2,6)+(0,2)=(6,20)$).}
  \item In one sentence: what does the $3$ in $L$ record?
        \ansline{The multiplier --- how many of row~1 we subtracted from row~2 during elimination.}
\end{enumerate}
\end{tcolorbox}

\begin{tcolorbox}[colback=white, colframe=black!40, title=\textbf{Tier A --- Approaching Proficiency},
    fonttitle=\bfseries, arc=1.5mm, left=3mm, right=3mm, top=2mm, bottom=2mm, breakable]
\textbf{One factorization, many orders.} A juice bar blends two mixes. Each liter of \textbf{Mix P}
carries $1$ unit sugar and $3$ units vitamin; each liter of \textbf{Mix Q} carries $2$ sugar and $8$
vitamin. With $p$ liters of P and $q$ of Q the totals are $A\begin{bsmallmatrix} p\\q \end{bsmallmatrix}=\bb$,
where $A=\begin{bmatrix} 1 & 2 \\ 3 & 8 \end{bmatrix}$ (sugar row, vitamin row).
\begin{enumerate}[leftmargin=1.7em, itemsep=6pt]
  \item Factor once: find $\ell$, and write $L$ and $U$.
        \ansline{$\ell=3$; $U=\begin{bsmallmatrix} 1&2\\0&2 \end{bsmallmatrix}$, $L=\begin{bsmallmatrix} 1&0\\3&1 \end{bsmallmatrix}$.}
  \item \textbf{Order 1:} sugar $4$, vitamin $14$. Two sweeps ($L\cc=\bb$, then $U\xx=\cc$) --- find $p,q$.
        \ansline{Forward: $c_1=4$, $3(4)+c_2=14\Rightarrow c_2=2$, so $\cc=\begin{bsmallmatrix} 4\\2 \end{bsmallmatrix}$. Back: $2q=2\Rightarrow q=1$, $p+2=4\Rightarrow p=2$. So $p=2$, $q=1$.}
  \item \textbf{Order 2:} sugar $7$, vitamin $27$. Same $L$ and $U$ --- find $p,q$.
        \ansline{Forward: $c_1=7$, $21+c_2=27\Rightarrow c_2=6$, so $\cc=\begin{bsmallmatrix} 7\\6 \end{bsmallmatrix}$. Back: $2q=6\Rightarrow q=3$, $p+6=7\Rightarrow p=1$. So $p=1$, $q=3$.}
  \item \emph{Interpret:} give each order's liters, and say why keeping $L$ and $U$ beats re-eliminating
        every order.
        \ansline{Order 1: $2$ L P, $1$ L Q; Order 2: $1$ L P, $3$ L Q. The recipe $A$ never changes, so
        one factorization $A=LU$ handles every order with just two quick sweeps --- no fresh elimination.}
\end{enumerate}
\end{tcolorbox}

\begin{tcolorbox}[colback=white, colframe=black!40, title=\textbf{Tier E --- Extension},
    fonttitle=\bfseries, arc=1.5mm, left=3mm, right=3mm, top=2mm, bottom=2mm, breakable]
\textbf{A $3\times3$ factorization.} Factor $A=\begin{bmatrix} 2 & 4 & 2 \\ 4 & 9 & 5 \\ 2 & 7 & 8 \end{bmatrix}$
by elimination, recording all three multipliers $\ell_{21},\ell_{31},\ell_{32}$ into $L$.
\ansline{$\ell_{21}=4/2=2$, $\ell_{31}=2/2=1$, then pivot $1$ gives $\ell_{32}=3/1=3$. So
$U=\begin{bsmallmatrix} 2&4&2\\0&1&1\\0&0&3 \end{bsmallmatrix}$, $L=\begin{bsmallmatrix} 1&0&0\\2&1&0\\1&3&1 \end{bsmallmatrix}$.}

\vspace{4pt}
\textbf{Then solve} $A\xx=\bb$ for $\bb=\begin{bmatrix} 10 \\ 22 \\ 19 \end{bmatrix}$ using two sweeps
($L\cc=\bb$ forward, then $U\xx=\cc$ back). \emph{Justify:} explain why writing down $L$ cost you no
extra arithmetic beyond the elimination you already did.
\ansline{Forward: $c_1=10$; $2(10)+c_2=22\Rightarrow c_2=2$; $10+3(2)+c_3=19\Rightarrow c_3=3$, so
$\cc=\begin{bsmallmatrix} 10\\2\\3 \end{bsmallmatrix}$. Back: $3z=3\Rightarrow z=1$; $y+z=2\Rightarrow y=1$;
$2x+4y+2z=10\Rightarrow x=2$. So $(x,y,z)=(2,1,1)$. $L$ is free because its entries are exactly the
multipliers you compute \emph{anyway} while eliminating --- you just write each one down instead of
discarding it.}
\end{tcolorbox}

\begin{teachernote}
Tier~A is the payoff: one $A=LU$, reused across right-hand sides --- the reason factoring beats
re-eliminating. Common slips: a minus sign in $L$ (it stores $+\ell$), or sweeping $U$ before $L$.
Tier~E makes the ``free $L$'' idea explicit --- the multipliers are a by-product of elimination, so
factoring adds no work. Tie back to 2.2: $L$ is the undo-matrices $E^{-1}$ collected together.
\end{teachernote}

\end{document}
